\chapter{基于物理感知深度强化学习的单目标自适应融合定位方法}
\label{ch:single_target_fusion}

\section{问题描述与总体框架}
\label{sec:3_1}

\subsection{单目标定位任务定义与场景约束}
\label{subsec:3_1_1}

本文研究的场景是：在复杂城市环境下，利用低成本无人机平台对地面非法低轨卫星终端进行非合作式无源单目标连续定位。与常规通信系统中的合作式定位不同，目标终端通常不提供同步辅助、位置信息或协议配合，无人机平台只能依赖接收到的目标相关多源观测信息对其位置进行估计。因此，本章所讨论的问题本质上属于复杂环境约束下的非合作式无源定位问题，其定位输出采用二维平面坐标形式，重点关注单目标连续轨迹的定位精度、轨迹平滑性以及对环境扰动的鲁棒性。

复杂城市环境中的高层建筑、狭窄街道以及大量反射界面会引起显著的非视距传播和多径效应，使得不同观测量的有效性随时间和空间持续变化。一方面，RSSI 易受局部遮挡、阴影衰落和快衰落影响，导致其与传播距离之间的映射关系失稳；另一方面，AoA 在强反射环境中容易偏离真实来波方向，而 PDR 虽然具有较好的短时平滑性，却会因惯性误差累积而在长时间尺度上逐渐偏离真实轨迹。上述现象表明，在复杂城市环境中，不同观测量之间并不存在固定不变的优劣关系，其可靠性会随着传播条件、遮挡状态和目标运动过程动态波动。因此，采用固定权重或静态融合规则难以适应复杂场景中的快速状态切换，这也是传统融合方法在实际排查任务中性能受限的重要原因。

基于上述分析，本章不再将单目标定位问题简单地视为多个观测量的线性组合问题，而是将其归结为一个面向动态环境的自适应融合决策问题。具体而言，本章关注的是：在复杂传播环境和多源异构观测条件下，如何根据信道状态和观测质量的实时变化，自适应地调整不同信息源的融合权重，从而获得稳定、连续且高精度的定位结果。围绕这一目标，本文后续将引入物理感知状态空间与深度强化学习决策机制，对单目标非法终端的动态融合定位进行研究。图~\ref{fig:single_target_scenario} 给出了本章所针对的单目标定位场景示意图。

\begin{figure}[htbp]
    \centering
    % \includegraphics[width=0.78\textwidth]{figures/ch3_single_target_scenario.pdf}
    \vspace{4.2cm}
    \caption{复杂城市环境下基于低成本无人机平台的单目标非法终端定位场景示意图}
    \label{fig:single_target_scenario}
\end{figure}

\subsection{多源异构定位信息加权融合模型}
\label{subsec:3_1_2}

为充分利用不同观测量之间的互补性，本文采用多源异构信息融合框架，将 RSSI、AoA 和 PDR 三类定位结果作为基础输入。之所以选择这三类观测，是因为它们分别从不同层面刻画目标位置特征：RSSI 反映信号功率层面的空间变化信息，AoA 提供几何方向约束，而 PDR 则利用惯性测量信息表征目标运动过程中的短时连续性。三者在信息来源和误差机理上具有明显差异，能够形成较强互补；但与此同时，其可靠性也会随着传播环境和运动状态变化而动态波动。因而，如何在不同场景下合理利用三类观测的优势，是实现稳定定位的关键。

为统一不同观测来源的输出形式，本文首先将各模块原始观测转换为对应的位置估计结果，再在位置层面进行加权融合。设在时刻 \(t\)，RSSI、AoA 和 PDR 模块分别给出位置估计 \(P_t^{\mathrm{RSSI}}\)、\(P_t^{\mathrm{AoA}}\) 和 \(P_t^{\mathrm{PDR}}\)，则最终融合结果可建模为
\begin{equation}
P_t^{\mathrm{Fused}}
=
\omega_t^{\mathrm{RSSI}} P_t^{\mathrm{RSSI}}
+
\omega_t^{\mathrm{AoA}} P_t^{\mathrm{AoA}}
+
\omega_t^{\mathrm{PDR}} P_t^{\mathrm{PDR}},
\label{eq:3_1_fusion_model}
\end{equation}
其中，\(\omega_t=[\omega_t^{\mathrm{RSSI}},\omega_t^{\mathrm{AoA}},\omega_t^{\mathrm{PDR}}]^T\) 为时刻 \(t\) 的权重向量，并满足归一化约束
\begin{equation}
\omega_t^{\mathrm{RSSI}}+\omega_t^{\mathrm{AoA}}+\omega_t^{\mathrm{PDR}}=1.
\label{eq:3_2_weight_constraint}
\end{equation}

上述模型的优势在于形式直观、实现简单，并且能够显式体现不同信息源在不同环境下的贡献差异。然而，该模型真正的难点并不在于融合形式本身，而在于权重 \(\omega_t\) 的动态确定机制。若采用固定权重，则模型难以适应复杂城市环境中瞬时变化的多径、遮挡和观测退化；若采用传统滤波器，则仍需依赖较强的噪声模型假设和先验统计条件，难以准确刻画非平稳环境中的观测质量波动。

进一步地，权重分配不仅依赖当前时刻不同观测量的质量，还与前一时刻的融合状态、轨迹连续性以及环境变化趋势相关。因此，该问题本质上并非一个静态参数估计问题，而是一个面向动态环境的序列决策问题。基于这一认识，本文不采用静态权重或人工规则，而是将权重生成过程交由智能体动态决策，并在下一小节中进一步将其建模为马尔可夫决策过程，为后续强化学习求解奠定基础。

\subsection{面向动态信道的马尔可夫决策过程建模}
\label{subsec:3_1_3}

由式~\eqref{eq:3_1_fusion_model} 可知，多源异构定位信息融合的关键在于权重向量 \(\omega_t\) 的动态确定。然而，在复杂城市环境中，这一权重分配过程并不是一个静态参数优化问题，而是一个具有显著时序性的动态决策问题。其原因在于：当前时刻的融合决策不仅依赖于当前环境状态和观测质量，还会进一步影响下一时刻的融合结果、误差传播以及轨迹连续性。换言之，在单目标连续定位过程中，系统的决策具有明显的“当前状态—当前动作—未来结果”链式关联，这使得多源融合过程天然具备序列决策特征。

具体而言，在复杂城市环境下，目标终端在运动过程中会不断经历视距（LOS）与非视距（NLOS）条件切换，RSSI、AoA 和 PDR 三类观测的可靠性也会随之动态变化。当智能体在某一时刻提高某一观测源的融合权重时，该动作不仅会改变当前时刻的融合位置输出，还会通过影响轨迹连续性、误差累积和后续状态演化，对下一时刻的决策产生作用。例如，在强反射环境下若仍赋予 AoA 过高权重，则会导致当前定位偏差增大，并使得下一时刻的几何一致性特征和轨迹误差进一步恶化；相反，若在 RSSI波动显著且 RMS 时延扩展快速升高时提前降低无线观测权重，则有助于抑制误差扩散并维持轨迹平滑性。因此，本章所研究的问题不是单步最优加权，而是在动态环境中根据状态演化持续更新最优策略的问题。

基于上述分析，本文将该动态融合过程抽象为马尔可夫决策过程（Markov Decision Process, MDP）。在这一框架下，智能体在每一时刻观测系统状态 \(S_t\)，并据此选择动作 \(a_t\)，动作执行后产生即时奖励 \(r_t\)，同时系统演化到下一状态 \(S_{t+1}\)。其中，状态用于刻画当前时刻的环境与观测条件，动作对应融合权重的更新策略，奖励则用于衡量当前决策在定位精度和轨迹平滑性方面的综合效果。只要当前决策所需的主要信息能够由当前状态向量近似表征，就可以将该问题视为一个近似满足马尔可夫性的序列决策问题。

为了形式化描述上述过程，本文将该问题表示为一个五元组
\begin{equation}
\mathcal{M}=\langle \mathcal{S},\mathcal{A},\mathcal{P},\mathcal{R},\gamma \rangle,
\label{eq:3_3_mdp_tuple}
\end{equation}
其中，\(\mathcal{S}\) 表示状态空间，\(\mathcal{A}\) 表示动作空间，\(\mathcal{P}\) 表示状态转移概率，\(\mathcal{R}\) 表示奖励函数，\(\gamma\in[0,1]\) 为折扣因子。对于本文而言，状态空间 \(\mathcal{S}\) 用于描述复杂传播环境下的物理信道状态、观测统计特征及融合历史信息；动作空间 \(\mathcal{A}\) 用于表示不同的权重更新策略；奖励函数 \(\mathcal{R}\) 则用于刻画当前融合决策对定位效果的优劣。

在该建模下，智能体的目标是学习一个最优策略 \(\pi^*\)，使得从任意状态出发时的长期累计回报最大化。对应的状态—动作价值函数可表示为
\begin{equation}
Q^{\pi}(s,a)
=
\mathbb{E}_{\pi}
\left[
\sum_{k=0}^{\infty}\gamma^k r_{t+k}
\,\middle|\,
S_t=s,\; a_t=a
\right].
\label{eq:3_4_q_definition}
\end{equation}
式~\eqref{eq:3_4_q_definition} 表明，当前动作的优劣不仅由即时奖励决定，还取决于该动作对未来一系列状态演化和奖励累积的长期影响。这一点与本文所研究的连续轨迹定位任务高度一致，因为单次权重更新往往会影响后续多个时刻的定位误差和轨迹平滑性。

依据 Bellman 最优性原理，最优状态—动作价值函数满足
\begin{equation}
Q^*(s,a)
=
\mathbb{E}
\left[
r_t+\gamma \max_{a'}Q^*(s',a')
\,\middle|\,
s,a
\right].
\label{eq:3_5_bellman}
\end{equation}
在实际求解中，本文采用深度 Q 网络对 \(Q^*(s,a)\) 进行近似。其核心思想是利用神经网络建立从状态到动作价值的映射，并通过经验回放和目标网络机制实现稳定训练。对应的参数更新损失可写为
\begin{equation}
L(\theta)
=
\mathbb{E}_{(s,a,r,s')\sim D}
\left[
\left(
y_i-Q(s,a;\theta)
\right)^2
\right],
\quad
y_i=r+\gamma\max_{a'}Q(s',a';\theta^-),
\label{eq:3_6_dqn_loss}
\end{equation}
其中，\(D\) 为经验回放池，\(\theta\) 为当前网络参数，\(\theta^-\) 为目标网络参数。

需要指出的是，本文并非将强化学习作为一般性的权重调节工具，而是将复杂城市环境下的多源定位融合过程视为一个随环境状态演化而不断更新的序列决策问题。在该问题中，传播条件具有明显的时变特性，不同观测量的可靠性会随遮挡、多径和目标运动过程发生变化，当前时刻的权重分配不仅决定当前定位结果，也会影响后续轨迹误差的传播特性。基于上述特点，采用马尔可夫决策过程对融合定位问题进行建模具有明确的合理性。在此基础上，后文将进一步分别构建 RL-IFA 基线算法与 PA-DQN 物理感知强化学习方法，并比较基于后验误差驱动的决策机制与引入物理先验后的决策机制在定位性能上的差异。图~\ref{fig:mdp_modeling_framework} 给出了本章单目标自适应融合定位问题的 MDP 建模示意图。

\begin{figure}[htbp]
    \centering
    % \includegraphics[width=0.82\textwidth]{figures/ch3_mdp_modeling_framework.pdf}
    \vspace{4.2cm}
    \caption{单目标自适应融合定位问题的马尔可夫决策过程建模示意图}
    \label{fig:mdp_modeling_framework}
\end{figure}

\section{复杂环境观测构造与物理先验信息提取}
\label{sec:3_2}

\subsection{射线追踪场景建模与环境先验获取}
\label{subsec:3_2_1}

为了建立观测波动与环境成因之间的对应关系，本文引入基于射线追踪的高保真场景建模方法，用于构造复杂城市环境中的传播先验信息。对于本章所研究的单目标连续定位任务而言，强化学习模型的训练与评估不仅需要大量可重复的交互样本，还要求在不同轨迹、不同信噪比以及不同视距/非视距切换条件下获得可控且带有真值标注的数据。然而，在真实城市环境中，大规模采集同时满足“多轨迹、可控环境切换、精确位置真值和多种信道状态”的连续定位数据代价较高，且实验可重复性有限。因此，本文采用“真实地理数据建模 + 射线追踪传播仿真”的方式构建复杂城市环境下的训练与测试场景，以兼顾物理一致性、样本规模和实验可控性。

在场景建模阶段，本文首先利用 OpenStreetMap（OSM）导出的真实地理信息数据，获取目标区域内建筑物的平面轮廓、高度信息及空间分布关系，并在 MATLAB Site Viewer 环境中完成三维城市几何场景重建。为了尽可能逼近真实传播条件，建模区域选取北京市典型建筑密集街区作为原型场景，通过将建筑物几何轮廓映射到三维坐标系中，形成包含高层建筑、街道走廊和开阔区域的城市峡谷环境。该建模过程为后续电磁传播仿真提供了几何基础，也使得不同位置的传播差异能够与真实城市结构建立对应关系。

\begin{figure}[htbp]
    \centering
    % \includegraphics[width=0.82\textwidth]{figures/ch3_osm_city_model.pdf}
    \vspace{4.2cm}
    \caption{基于 OSM 数据的复杂城市三维场景建模示意图}
    \label{fig:ch3_osm_city_model}
\end{figure}

\subsection{多源观测量生成与特征表达}
\label{subsec:3_2_2}

在本章所构建的自适应融合定位框架中，系统并不直接以原始射频观测作为最终输入，而是将不同来源的观测信息分别转化为可用于定位的初步估计结果，并在此基础上完成后续的动态加权融合。关于 RSSI、AoA 以及 PDR 等观测量的获取方式、基本测量机理和定位模型，前文第二章已经进行了讨论。因此，本节不再重复展开各类观测的底层测量过程，而是在前述理论基础上，重点说明 RSSI、AoA 和 PDR 三类异构观测在本章中的生成形式、表达方式及其在融合框架中的作用。

结合当前研究基础，本文选取 RSSI、AoA 和 PDR 三类观测作为主要输入源。其原因在于：RSSI 能够反映信号功率层面的空间变化特征，AoA 提供几何方向约束，而 PDR 则利用惯性信息表征目标短时运动连续性。三者分别对应功率域、几何域和运动域信息，具有较强互补性，能够共同描述复杂城市环境下单目标定位所需的多维特征。然而，这三类观测在可靠性上并不存在固定不变的优劣关系，而是会随着传播条件、遮挡状态和目标运动过程持续变化，因此需要在统一框架下对其进行动态融合。

需要指出的是，本文对三类观测的使用包含两条并行数据链路。第一条为定位输出链，即 RSSI、AoA 和 PDR 模块分别根据各自观测模型输出初步位置估计 \(P_t^{\mathrm{RSSI}}\)、\(P_t^{\mathrm{AoA}}\) 和 \(P_t^{\mathrm{PDR}}\)，并作为后续加权融合模块的直接输入。第二条为状态构造链，即各模块在生成初步位置估计的同时，保留其原始观测数据或中间统计量，并与上一节得到的环境先验共同送入物理感知特征提取器，用于构造强化学习智能体的状态向量。也就是说，本节中的多源观测不仅服务于位置解算，还承担了状态感知输入的功能，这是本章区别于传统静态融合方法的重要特点。

\begin{figure}[htbp]
    \centering
    % \includegraphics[width=0.84\textwidth]{figures/ch3_multi_source_observation_flow.pdf}
    \vspace{4.2cm}
    \caption{多源观测生成与融合输入构造流程示意图}
    \label{fig:multi_source_observation_flow}
\end{figure}

\begin{figure}[htbp]
    \centering
    % \includegraphics[width=0.84\textwidth]{figures/ch3_parallel_observation_state_chain.pdf}
    \vspace{4.2cm}
    \caption{定位输出链与状态构造链的并行结构关系示意图}
    \label{fig:multi_source_parallel_chain}
\end{figure}

\subsection{RMS 时延扩展与几何一致性特征构造}
\label{subsec:3_2_3}

为了使强化学习智能体能够在定位误差显著累积之前感知环境变化，本文并未仅依赖后验定位误差或当前融合结果构造状态空间，而是进一步引入能够反映传播机理、观测稳定性和系统几何一致性的多层状态特征。基于这一认识，本文构建由物理层特征、统计层特征、几何层特征以及历史决策信息共同组成的分层状态表示。

首先，本文选取 RMS 时延扩展 \(\sigma_\tau\) 作为表征多径复杂度的核心物理层特征。该特征来源于上一节射线追踪场景建模得到的路径级传播参数，能够反映接收信号在时延域上的离散程度。RMS 时延扩展定义为
\begin{equation}
\sigma_\tau=
\sqrt{
\frac{\sum_{n=1}^{N} |\alpha_n|^2 (\tau_n-\bar{\tau})^2}
{\sum_{n=1}^{N} |\alpha_n|^2}
},
\label{eq:3_8_rms_delay}
\end{equation}
其中，\(N\) 为有效多径分量数，\(\alpha_n\) 和 \(\tau_n\) 分别表示第 \(n\) 条路径的幅度和绝对时延，\(\bar{\tau}\) 为功率加权平均时延。

其次，为刻画接收功率在时间上的波动程度，本文引入 RSSI 方差 \(\sigma^2_{\mathrm{RSSI}}\) 作为统计层特征。设在长度为 \(W\) 的滑动窗口内，接收功率序列为 \(\{RSSI_k\}\)，则其方差定义为
\begin{equation}
\sigma^2_{\mathrm{RSSI}}
=
\frac{1}{W}
\sum_{k=t-W+1}^{t}
\left(
RSSI_k-\overline{RSSI}_t
\right)^2,
\label{eq:3_9_rssi_var}
\end{equation}
其中，\(\overline{RSSI}_t\) 为窗口内平均接收功率。

进一步地，为利用轨迹连续性对无线观测结果进行系统层约束，本文构造几何一致性特征 \(D_{\mathrm{RF-PDR}}\)。其中，\(P_{\mathrm{RF}}\) 表示当前时刻由无线电侧得到的位置估计结果，\(P_{\mathrm{PDR}}\) 表示由惯性递推获得的 PDR 位置估计。两者之间的欧氏距离定义为
\begin{equation}
D_{\mathrm{RF-PDR}}
=
\|P_{\mathrm{RF}}-P_{\mathrm{PDR}}\|_2
=
\sqrt{
(x_{\mathrm{RF}}-x_{\mathrm{PDR}})^2
+
(y_{\mathrm{RF}}-y_{\mathrm{PDR}})^2
},
\label{eq:3_10_geo_consistency}
\end{equation}

在上述基础上，本文将物理层特征、统计层特征、几何层特征与上一时刻权重向量共同组合，形成当前时刻的物理感知状态向量
\begin{equation}
S_t=
\left[
\sigma_\tau,\,
\sigma^2_{\mathrm{RSSI}},\,
D_{\mathrm{RF-PDR}},\,
\omega_{t-1}
\right]^T,
\label{eq:3_11_state_vector}
\end{equation}
其中，\(\omega_{t-1}\) 表示上一时刻的融合权重。图~\ref{fig:ch3_physics_aware_state_space} 给出了物理感知状态空间构成示意图。

\begin{figure}[htbp]
    \centering
    % \includegraphics[width=0.84\textwidth]{figures/ch3_physics_aware_state_space.pdf}
    \vspace{4.2cm}
    \caption{物理感知状态空间构成示意图}
    \label{fig:ch3_physics_aware_state_space}
\end{figure}

\section{基于 Q-Learning 的 RL-IFA 基线自适应融合算法}
\label{sec:3_3}

\subsection{状态空间设计}
\label{subsec:3_3_1}

为考察物理感知状态构造对融合定位性能的影响，本文首先构建一种不显式引入传播环境先验的强化学习信息融合算法，记为基于强化学习的信息融合算法（Reinforcement Learning-based Information Fusion Algorithm, RL-IFA），并将其作为后续 PA-DQN 方法的基线模型。RL-IFA 的基本思想是：在多源观测融合过程中，将当前融合结果的定位误差作为状态反馈量，利用 Q-Learning 学习不同观测量之间的权重调整策略，从而实现自适应融合定位。

与后续 PA-DQN 的物理感知状态空间不同， RL-IFA 并不对多径结构、观测稳定性或几何一致性进行显式建模，而仅依据当前融合结果与真实位置之间的偏差来刻画系统状态。设时刻 \(t\) 的融合定位结果为
\(
P_t^{\mathrm{Fused}}=(x_t^{\mathrm{Fused}},y_t^{\mathrm{Fused}})
\)，
对应真实位置为
\(
P_t^{\mathrm{GT}}=(x_t^{\mathrm{GT}},y_t^{\mathrm{GT}})
\)，
则当前时刻的融合误差定义为
\begin{equation}
\varepsilon_t
=
\left\|P_t^{\mathrm{Fused}}-P_t^{\mathrm{GT}}\right\|_2
=
\sqrt{
\left(x_t^{\mathrm{GT}}-x_t^{\mathrm{Fused}}\right)^2
+
\left(y_t^{\mathrm{GT}}-y_t^{\mathrm{Fused}}\right)^2
}.
\label{eq:3_12_rlifa_error}
\end{equation}

式~\eqref{eq:3_12_rlifa_error} 所定义的误差量直接反映了当前融合策略下的定位效果，因此可作为强化学习决策的后验反馈依据。考虑到经典 Q-Learning 通常在离散状态空间中实现，为控制状态空间规模并提高训练可行性，本文对连续误差 \(\varepsilon_t\) 进行离散化处理。参考相关研究中对融合误差状态化的处理方式，本文将误差区间 \([0,1)\,\mathrm{m}\) 均匀划分为 100 个离散状态，并将大于等于 1 m 的误差统一映射为终端状态。于是，时刻 \(t\) 的状态 \(s_t\) 定义为
\begin{equation}
s_t=
\begin{cases}
\mathrm{round}(100\,\varepsilon_t), & 0 \leq \varepsilon_t < 1,\\
100, & \varepsilon_t \geq 1.
\end{cases}
\label{eq:3_13_rlifa_state}
\end{equation}

上述状态定义方式具有两方面特点。其一，状态的构造仅依赖融合结果与真实位置之间的后验误差，因而实现简单，便于形成标准 Q-Learning 基线；其二，状态编号随误差大小单调变化，从而能够较直接地反映当前融合效果的优劣。图~\ref{fig:rlifa_state_design} 给出了 RL-IFA 基线算法中基于融合误差的状态空间设计示意图。

\begin{figure}[htbp]
    \centering
    % \includegraphics[width=0.80\textwidth]{figures/ch3_rlifa_state_design.pdf}
    \vspace{4.2cm}
    \caption{RL-IFA 基线算法中基于融合误差的状态空间设计示意图}
    \label{fig:rlifa_state_design}
\end{figure}

\subsection{动作空间设计}
\label{subsec:3_3_2}

在 RL-IFA 基线算法中，动作空间的作用是根据当前状态对多源观测的融合权重进行更新。由于本文采用 Q-Learning 进行求解，为保证动作集合有限且便于训练实现，本文对权重调整过程进行离散化处理。具体而言，动作并不直接给出新的连续权重值，而是规定在当前时刻对各观测源权重执行“增加”“减小”或“保持不变”三类基本操作。通过这种方式，多源融合权重优化问题可被转化为有限动作空间下的离散决策问题。

考虑到三类观测的融合权重满足归一化约束
\begin{equation}
\omega_t^{\mathrm{RSSI}}+\omega_t^{\mathrm{AoA}}+\omega_t^{\mathrm{PDR}}=1,
\label{eq:3_14_rlifa_weight_constraint}
\end{equation}
因此，在动作设计中无须对三个权重同时独立更新。为降低动作空间维数并避免冗余组合，本文仅对 \(\omega_t^{\mathrm{RSSI}}\) 和 \(\omega_t^{\mathrm{AoA}}\) 施加离散动作，第三个权重 \(\omega_t^{\mathrm{PDR}}\) 由归一化约束自动确定。

设权重更新步长参数为 \(\kappa\)，则对任一可调权重 \(\omega_t^{x}\)（其中 \(x \in \{\mathrm{RSSI},\mathrm{AoA}\}\)），其更新规则定义为
\begin{equation}
\omega_{t+1}^{x}=
\begin{cases}
(1+\kappa)\omega_t^{x}, & \text{若执行增加动作},\\[4pt]
(1-\kappa)\omega_t^{x}, & \text{若执行减小动作},\\[4pt]
\omega_t^{x}, & \text{若执行保持动作}.
\end{cases}
\label{eq:3_15_rlifa_weight_update}
\end{equation}
在得到 \(\omega_{t+1}^{\mathrm{RSSI}}\) 和 \(\omega_{t+1}^{\mathrm{AoA}}\) 后，由式~\eqref{eq:3_14_rlifa_weight_constraint} 可进一步计算
\begin{equation}
\omega_{t+1}^{\mathrm{PDR}}
=
1-\omega_{t+1}^{\mathrm{RSSI}}-\omega_{t+1}^{\mathrm{AoA}}.
\label{eq:3_16_rlifa_pdr_weight}
\end{equation}

基于上述定义，本文将 RL-IFA 的动作空间构造为 9 个离散动作，即对 \(\omega^{\mathrm{RSSI}}\) 和 \(\omega^{\mathrm{AoA}}\) 分别执行“增加 / 减小 / 保持”三类操作后形成的组合集合：
\begin{equation}
\mathcal{A}
=
\left\{
a_1,a_2,\ldots,a_9
\right\}.
\label{eq:3_17_rlifa_action_set}
\end{equation}
这 9 个动作的具体定义如表~\ref{tab:rlifa_action_space} 所示。由表可见，动作空间覆盖了两路权重的全部基本组合情况，包括同步增强、同步抑制以及一增一减、一增一稳、一减一稳等多种更新模式。

\begin{table}[htbp]
    \centering
    \caption{RL-IFA 基线算法中离散动作空间定义}
    \label{tab:rlifa_action_space}
    \begin{tabular}{cccc}
        \toprule
        \textbf{动作编号} & \(\boldsymbol{\omega^{\mathrm{RSSI}}}\) \textbf{更新} & \(\boldsymbol{\omega^{\mathrm{AoA}}}\) \textbf{更新} & \(\boldsymbol{\omega^{\mathrm{PDR}}}\) \textbf{更新方式} \\
        \midrule
        \(a_1\) & 增加 & 增加 & 由归一化约束计算 \\
        \(a_2\) & 增加 & 减小 & 由归一化约束计算 \\
        \(a_3\) & 增加 & 保持 & 由归一化约束计算 \\
        \(a_4\) & 减小 & 增加 & 由归一化约束计算 \\
        \(a_5\) & 减小 & 减小 & 由归一化约束计算 \\
        \(a_6\) & 减小 & 保持 & 由归一化约束计算 \\
        \(a_7\) & 保持 & 增加 & 由归一化约束计算 \\
        \(a_8\) & 保持 & 减小 & 由归一化约束计算 \\
        \(a_9\) & 保持 & 保持 & 由归一化约束计算 \\
        \bottomrule
    \end{tabular}
\end{table}

图~\ref{fig:rlifa_action_design} 给出了 RL-IFA 基线算法中动作空间与权重更新关系的示意图。

\begin{figure}[htbp]
    \centering
    % \includegraphics[width=0.82\textwidth]{figures/ch3_rlifa_action_design.pdf}
    \vspace{4.2cm}
    \caption{RL-IFA 基线算法中动作空间与权重更新关系示意图}
    \label{fig:rlifa_action_design}
\end{figure}

\subsection{奖励函数设计}
\label{subsec:3_3_3}

在 RL-IFA 中，状态用于描述当前融合结果的误差水平，动作用于调整各观测分支的权重分配，而奖励函数则负责评价该次权重调整是否有利于定位结果改善。由于基线算法的目标是验证强化学习在多源异构观测融合中的有效性，因此本节奖励设计首先围绕定位精度展开，即将融合输出与真实位置之间的偏差作为策略优劣的直接判据。

设第 \(k\) 个时刻获取的三类定位观测分支经加权融合后计算得到的位置估计为
\(\hat{\mathbf{p}}_k^{\mathrm{Fused}}=[\hat{x}_k^{\mathrm{Fused}},\hat{y}_k^{\mathrm{Fused}}]^T\)，
对应真实位置为
\(\mathbf{p}_k^{\mathrm{GT}}=[x_k^{\mathrm{GT}},y_k^{\mathrm{GT}}]^T\)，
则融合定位误差可写为
\begin{equation}
e_k=\left\|\hat{\mathbf{p}}_k^{\mathrm{Fused}}-\mathbf{p}_k^{\mathrm{GT}}\right\|_2
=\sqrt{\left(\hat{x}_k^{\mathrm{Fused}}-x_k^{\mathrm{GT}}\right)^2+
\left(\hat{y}_k^{\mathrm{Fused}}-y_k^{\mathrm{GT}}\right)^2 }.
\label{eq:rlifa_reward_error}
\end{equation}

考虑到 Q-Learning 采用表格式价值迭代进行策略更新，奖励函数不宜设计得过于复杂。本文采用与定位误差倒数相关的分段奖励形式：
\begin{equation}
r_k=
\begin{cases}
100, & e_k=0,\\[4pt]
\mathrm{round}\!\left(\dfrac{100}{e_k}\right), & 0<e_k<1,\\[8pt]
-100, & e_k\geq 1.
\end{cases}
\label{eq:rlifa_reward}
\end{equation}

式~\eqref{eq:rlifa_reward} 的设计与前述状态离散方式保持一致。当融合误差较小时，奖励值取正，且误差越小，奖励越大；当误差达到或超过 \(1\,\mathrm{m}\) 后，给予固定负奖励。

\subsection{权重更新机制}
\label{subsec:3_3_4}

在完成状态空间、动作空间与奖励函数定义后，RL-IFA 的核心问题可进一步归结为：智能体如何依据当前误差状态选择合适动作，并将该动作映射为对融合权重的实际更新，从而在连续交互过程中逐步学习得到更优的加权融合策略。与固定权重方法不同，RL-IFA 并不预先假设 RSSI、AoA 和 PDR 三类观测在全局范围内具有恒定的重要性，而是将其权重分配过程建模为一个随时间滚动更新的离散序列决策过程。对于时刻 \(t\)，智能体首先根据当前融合误差对应的离散状态 \(s_t\) 查询动作价值，再执行相应的权重调整动作，随后由新的融合结果产生下一时刻状态与即时奖励，并进一步用于 Q 表更新。由此，整个系统形成“状态观测—动作选择—权重更新—结果反馈—价值迭代”的闭环学习机制。

具体而言，设上一时刻的融合权重向量为
\begin{equation}
\boldsymbol{\omega}_t
=
\left[
\omega_t^{\mathrm{RSSI}},
\omega_t^{\mathrm{AoA}},
\omega_t^{\mathrm{PDR}}
\right]^T,
\label{eq:3_19_rlifa_weight_vector}
\end{equation}
其中各分量满足非负性与归一化约束。根据前文 3.3.2 中的动作空间设计，智能体在状态 \(s_t\) 下从动作集合
\(\mathcal{A}=\{a_1,a_2,\ldots,a_9\}\)
中选取一个动作 \(a_t\)，并对
\(\omega_t^{\mathrm{RSSI}}\)
和
\(\omega_t^{\mathrm{AoA}}\)
实施“增加 / 减小 / 保持”三类基本操作，再由归一化约束计算
\(\omega_t^{\mathrm{PDR}}\)。
因此，RL-IFA 的权重更新并不是连续优化求解，而是通过有限动作集合对当前权重向量进行离散修正。

设比例更新步长为 \(\kappa\)，对于可调权重分量
\(\omega_t^{x}\)（其中 \(x\in\{\mathrm{RSSI},\mathrm{AoA}\}\)），在动作 \(a_t\) 的作用下，其更新结果可写为
\begin{equation}
\omega_{t+1}^{x}=
\begin{cases}
(1+\kappa)\omega_t^{x}, & \text{若动作规定该分量增加},\\[4pt]
(1-\kappa)\omega_t^{x}, & \text{若动作规定该分量减小},\\[4pt]
\omega_t^{x}, & \text{若动作规定该分量保持不变}.
\end{cases}
\label{eq:3_20_rlifa_weight_update}
\end{equation}
在得到
\(\omega_{t+1}^{\mathrm{RSSI}}\)
和
\(\omega_{t+1}^{\mathrm{AoA}}\)
之后，第三个分量由归一化关系确定为
\begin{equation}
\omega_{t+1}^{\mathrm{PDR}}
=
1-\omega_{t+1}^{\mathrm{RSSI}}-\omega_{t+1}^{\mathrm{AoA}}.
\label{eq:3_21_rlifa_weight_pdr}
\end{equation}

将更新后的权重代入位置级融合模型，可得下一时刻的融合位置输出
\begin{equation}
P_{t+1}^{\mathrm{Fused}}
=
\omega_{t+1}^{\mathrm{RSSI}} P_{t+1}^{\mathrm{RSSI}}
+
\omega_{t+1}^{\mathrm{AoA}} P_{t+1}^{\mathrm{AoA}}
+
\omega_{t+1}^{\mathrm{PDR}} P_{t+1}^{\mathrm{PDR}}.
\label{eq:3_22_rlifa_fused_next}
\end{equation}
由式~\eqref{eq:3_22_rlifa_fused_next} 进一步计算融合误差
\(\varepsilon_{t+1}\)，即可依据 3.3.1 中的离散映射规则获得下一状态 \(s_{t+1}\)，并依据 3.3.3 中定义的奖励函数获得即时奖励 \(r_t\)。由此，动作 \(a_t\) 对系统的影响被完整地传递到“新状态—新奖励”这一强化学习更新环节中。

在策略执行阶段，RL-IFA 采用 \(\varepsilon\)-greedy 机制在“探索”与“利用”之间进行平衡。具体来说，在状态 \(s_t\) 下，智能体以概率 \(\varepsilon\) 从动作集合 \(\mathcal{A}\) 中随机选择一个动作，以概率 \(1-\varepsilon\) 选择当前 Q 值最大的动作，即
\begin{equation}
a_t=
\begin{cases}
\text{random action}, & \text{with probability } \varepsilon,\\[4pt]
\arg\max\limits_{a\in\mathcal{A}}Q(s_t,a), & \text{with probability } 1-\varepsilon.
\end{cases}
\label{eq:3_23_rlifa_epsilon_greedy}
\end{equation}
引入该策略的原因在于，若始终选择当前最优动作，则智能体在训练初期容易陷入局部最优的权重配置；而适当保留随机探索能力，则有助于其在不同误差状态下遍历更多的动作组合，从而更充分地学习各观测源在不同条件下的相对贡献。

在获得动作执行后的即时奖励 \(r_t\) 与下一状态 \(s_{t+1}\) 后，RL-IFA 采用标准 Q-Learning 规则对状态—动作价值函数进行迭代更新。设 \(Q(s_t,a_t)\) 表示在状态 \(s_t\) 下执行动作 \(a_t\) 的价值估计，则其更新公式为
\begin{equation}
Q(s_t,a_t)
\leftarrow
Q(s_t,a_t)
+
\alpha
\left[
r_t
+
\gamma \max_{a\in\mathcal{A}}Q(s_{t+1},a)
-
Q(s_t,a_t)
\right],
\label{eq:3_24_rlifa_q_update}
\end{equation}
其中，\(\alpha\) 为学习率，\(\gamma\) 为折扣因子。式~\eqref{eq:3_24_rlifa_q_update} 表明，RL-IFA 对当前动作优劣的判断并不只取决于其带来的单步误差变化，还取决于执行该动作后系统未来可能获得的长期累计回报。也就是说，智能体在更新权重时并不是简单地“哪个动作让当前误差更小就选哪个”，而是在多步交互意义下逐步逼近能够带来更优长期融合效果的动作选择策略。

从算法流程上看，RL-IFA 的单步权重更新可概括为如下过程：首先，根据当前融合结果与真实位置之间的误差 \(\varepsilon_t\)，由式~\eqref{eq:3_13_rlifa_state} 将其映射为离散状态 \(s_t\)；其次，基于当前 Q 表和 \(\varepsilon\)-greedy 策略，从动作空间中选择动作 \(a_t\)；随后，依据式~\eqref{eq:3_20_rlifa_weight_update} 和式~\eqref{eq:3_21_rlifa_weight_pdr} 对权重向量进行更新，并计算新的融合位置 \(P_{t+1}^{\mathrm{Fused}}\)；接着，根据新的融合误差生成奖励 \(r_t\) 和下一状态 \(s_{t+1}\)；最后，利用式~\eqref{eq:3_24_rlifa_q_update} 对 Q 表进行迭代修正，并进入下一轮决策。随着轨迹样本不断积累，Q 表中不同状态—动作对的价值差异将逐步拉开，最终形成从误差状态到权重更新动作的经验映射关系。

需要指出的是，RL-IFA 的权重更新机制本质上仍然是一种基于后验误差反馈的自适应调整过程。其优点在于结构清晰、实现简单，且能够在不依赖显式传播模型的前提下，通过试错学习获得比固定权重更灵活的融合策略。对于作为基线算法的 RL-IFA 而言，这一机制已经足以验证强化学习在多源异构观测自适应融合中的可行性。然而，其局限也同样明显：由于状态仅由融合误差离散化得到，动作决策的依据仍然停留在结果层面，智能体只能在误差已经发生之后再通过权重更新进行补偿，而无法直接感知当前环境中多径增强、遮挡变化或观测退化的物理成因。因此，该权重更新机制虽然能够实现“误差驱动”的动态调权，但尚不具备“环境感知”的前瞻决策能力。

此外，RL-IFA 的动作更新采用比例缩放方式，虽然能够避免权重完全固定，但其控制对象仍然是当前时刻的绝对权重值，而不是相邻时刻之间的平滑增量。这意味着当误差状态在若干离散区间之间频繁跳变时，智能体可能在不同动作之间反复切换，从而导致融合权重产生较明显的震荡。这种震荡在单步误差上未必总是表现为显著恶化，但在连续轨迹层面往往会引起局部波动甚至锯齿状偏折。也正因如此，RL-IFA 更适合作为后续改进方法的基线模型：它建立了“状态—动作—奖励—权重更新”的基本强化学习融合框架，但尚未解决复杂城市环境下连续定位任务所要求的环境感知能力与轨迹平滑性约束问题。

综上，RL-IFA 通过 Q-Learning 实现了对融合权重的闭环在线更新，其核心机制在于：将当前后验误差映射为离散状态，通过 \(\varepsilon\)-greedy 选择调权动作，并依据即时奖励与下一状态对 Q 表进行价值迭代，从而逐步学习得到不同误差等级下的最优权重更新策略。该机制为后续 PA-DQN 的改进提供了统一基线。与 RL-IFA 仅依赖后验误差反馈不同，下一节将在此基础上进一步引入物理感知状态空间、增量式动作空间与联合奖励设计，以实现面向复杂城市环境的更稳定、更具前瞻性的自适应融合决策。

\section{融合物理先验知识的 PA-DQN 智能体设计}

在前述小节中，本文已经分别给出了 PA-DQN 的物理感知状态空间、增量式动作空间以及奖励函数设计。至此，单目标自适应融合定位问题的序列决策描述已经基本完成。在前述小节中，智能体在每一时刻根据环境感知状态 \(\mathbf{S}_t\) 选择权重调整动作 \(a_t\)，并依据奖励反馈不断修正决策策略。进一步需要解决的问题在于，如何在高维连续状态空间条件下有效逼近状态--动作价值函数，并据此实现对最优融合策略的学习。

对于 RL-IFA 基线算法而言，其状态和动作均经过离散化处理，因此可以采用表格式 Q-Learning 直接对状态--动作值函数进行存储与更新。然而，PA-DQN 中的状态空间已不再是单一误差驱动的低维离散状态，而是扩展为包含物理层、统计层、几何层和历史决策信息的连续高维向量。若仍采用表格式方法，不仅会导致状态划分过粗、信息表达不足，还会使状态规模随特征维数增加而迅速膨胀，无法满足复杂城市环境下连续定位任务的求解需求。因此，本文采用深度 Q 网络（Deep Q-Network, DQN）对状态--动作价值函数进行近似表示，在 RL-IFA 序列决策框架基础上构建物理感知深度 Q 网络模型。
\subsection{物理感知状态空间设计}



\label{subsec:3_4_1}





基于上述考虑，本文将第 \(t\) 个时刻的状态向量定义为

\begin{equation}
\mathbf{S}_t=
\left[
\sigma_{\tau,t},\,
\sigma^2_{\mathrm{RSSI},t},\,
D_{\mathrm{RF\mbox{-}PDR},t},\,
\boldsymbol{\omega}_{t-1}
\right]^T,
\label{eq:pa_dqn_state}
\end{equation}

其中，\(\sigma_{\tau,t}\) 为 RMS 时延扩展，\(\sigma^2_{\mathrm{RSSI},t}\) 为 RSSI 方差，\(D_{\mathrm{RF\mbox{-}PDR},t}\) 为无线观测定位结果与轨迹推算结果之间的几何一致性特征，\(\boldsymbol{\omega}_{t-1}\) 为上一时刻的融合权重向量。上述各特征的具体构造方法已在 3.2.3 中给出，此处不再重复展开。



式~\eqref{eq:pa_dqn_state} 的关键不在于状态维数增加，而在于状态信息由“结果误差”转向“环境表征”。其中，\(\sigma_{\tau,t}\) 反映当前传播路径在时延域上的离散程度。当多径成分增多、主导路径不再稳定时，该特征通常会增大，对应场景往往偏离理想视距传播条件。将其引入状态后，智能体可以在误差尚未明显扩大之前，先对当前传播环境的复杂程度作出响应。



\(\sigma^2_{\mathrm{RSSI},t}\) 用于表征接收信号强度在局部时间窗口内的波动情况。对于 RSSI 分支而言，绝对功率水平只能部分反映距离信息，而方差更能体现当前观测的稳定程度。当遮挡、深衰落或局部散射增强时，RSSI 序列通常会出现更明显的起伏，此时若仍维持较高的 RSSI 融合权重，往往会降低最终定位结果的稳定性。因此，RSSI 方差进入状态空间后，智能体能够对功率观测的可信度作出更细致的区分。



\(D_{\mathrm{RF\mbox{-}PDR},t}\) 则用于描述不同定位分支之间的几何协调关系。与单独观察某一路观测质量不同，该特征反映的是无线观测侧位置估计与轨迹推算结果之间的一致程度。当两者偏差较小时，说明当前观测结果与轨迹连续性约束基本一致；反之，若二者差异明显，则意味着至少有一类定位分支已偏离当前运动状态。将这一特征纳入状态后，智能体不仅能感知单路观测是否稳定，还能进一步判断不同信息源之间是否相互支持。



除环境和观测特征外，本文还将上一时刻的权重向量 \(\boldsymbol{\omega}_{t-1}\) 作为状态的一部分。连续定位过程中，相邻时刻的观测条件通常具有相关性，权重更新也不应完全脱离上一时刻的决策结果。若状态中不包含历史权重信息，则当前动作只能依据瞬时特征独立决策，容易导致相邻时刻权重出现大幅跳变。引入 \(\boldsymbol{\omega}_{t-1}\) 后，智能体在进行当前决策时能够保留必要的历史上下文，这也为后续增量式动作设计提供了基础。



从方法层看，PA-DQN 的状态空间重构主要体现为两个变化。其一，状态输入由 RL-IFA 中的单一后验误差，扩展为传播、统计、几何和历史决策四类信息的联合描述；其二，状态的作用由“反映当前结果好坏”转为“支撑当前权重应如何调整”。前者决定了智能体能够获得更丰富的环境信息，后者决定了这些信息可以直接服务于后续策略学习。这样处理后，权重分配不再只是对已发生误差的被动修正，而更接近于依据当前环境特征进行前置调节。



需要说明的是，本节给出的状态空间并未改变位置级动态加权融合模型本身，变化仅发生在决策层。RSSI、AoA 和 PDR 三路位置估计仍按既定方式输出初始定位结果，PA-DQN 只是利用式~\eqref{eq:pa_dqn_state} 所构造的状态向量对当前环境进行表征，并在此基础上学习更合适的权重更新策略。后续将在此基础上进一步给出 PA-DQN 的动作空间设计与奖励函数构造。




\subsection{物理感知状态空间设计}

\label{subsec:3_4_1}

RL-IFA 将融合定位误差离散化后作为智能体状态输入，其优点在于建模形式简单，便于采用表格式 Q-Learning 完成权重学习。但对于复杂城市环境下的单目标连续定位任务，这种状态定义仍然存在明显局限。首先，误差状态属于后验反馈，只有当融合结果已经偏离真实位置后，智能体才会通过奖励变化感知当前权重配置存在问题。其次，不同环境条件下产生的定位误差在数值上可能相近，但其成因并不一致。例如，RSSI 波动增强、AoA 主径偏置和 PDR 累积漂移都可能导致误差上升，仅凭误差大小本身，智能体无法区分当前失真主要来源于哪一路观测。这样得到的策略更接近于事后修正，而不是面向环境变化的主动调节。
为提高权重决策对复杂传播条件的适应性，我们设计了PA-DQN。PA-DQN 不再沿用单一误差驱动的状态定义，而是在状态层显式引入能够反映传播复杂度、信号稳定性、跨模态一致性以及历史决策信息的物理感知特征。这里的“物理感知”并不是重新构造新的定位观测量，而是将前文 3.2.3 已提取的环境相关特征引入决策过程，使智能体在更新融合权重之前，先对当前观测质量和环境状态形成更有针对性的描述。

基于上述考虑，本文将第 \(t\) 个时刻的状态向量定义为
\begin{equation}
\mathbf{S}_t=
\left[
\sigma_{\tau,t},\,
\sigma^2_{\mathrm{RSSI},t},\,
D_{\mathrm{RF\mbox{-}PDR},t},\,
\boldsymbol{\omega}_{t-1}
\right]^T,
\label{eq:pa_dqn_state}
\end{equation}
其中，\(\sigma_{\tau,t}\) 为 RMS 时延扩展，\(\sigma^2_{\mathrm{RSSI},t}\) 为 RSSI 方差，\(D_{\mathrm{RF\mbox{-}PDR},t}\) 为无线观测定位结果与轨迹推算结果之间的几何一致性特征，\(\boldsymbol{\omega}_{t-1}\) 为上一时刻的融合权重向量。上述各特征的具体构造方法已在 3.2.3 中给出，此处不再重复展开。

式~\eqref{eq:pa_dqn_state} 不仅状态维数增加，还是的状态信息由结果误差转向环境物理信息表征。其中，\(\sigma_{\tau,t}\) 反映当前传播路径在时延域上的离散程度。当多径成分增多、主导路径不再稳定时，该特征通常会增大，对应场景往往偏离理想视距传播条件。将其引入状态后，智能体可以在误差尚未明显扩大之前，先对当前传播环境的复杂程度作出响应。

\(\sigma^2_{\mathrm{RSSI},t}\) 用于表征接收信号强度在局部时间窗口内的波动情况。对于 RSSI 分支而言，绝对功率水平只能部分反映距离信息，而方差更能体现当前观测的稳定程度。当遮挡、深衰落或局部散射增强时，RSSI 序列通常会出现更明显的起伏，此时若仍维持较高的 RSSI 融合权重，往往会降低最终定位结果的稳定性。因此，RSSI 方差进入状态空间后，智能体能够对功率观测的可信度作出更细致的区分。

\(D_{\mathrm{RF\mbox{-}PDR},t}\) 则用于描述不同定位分支之间的几何协调关系。与单独观察某一路观测质量不同，该特征反映的是无线观测侧位置估计与轨迹推算结果之间的一致程度。当两者偏差较小时，说明当前观测结果与轨迹连续性约束基本一致；反之，若二者差异明显，则意味着至少有一类定位分支已偏离当前运动状态。将这一特征纳入状态后，智能体不仅能感知单路观测是否稳定，还能进一步判断不同信息源之间是否相互支持。

除环境和观测特征外，本文还将上一时刻的权重向量 \(\boldsymbol{\omega}_{t-1}\) 作为状态的一部分。连续定位过程中，相邻时刻的观测条件通常具有相关性，权重更新也不应完全脱离上一时刻的决策结果。若状态中不包含历史权重信息，则当前动作只能依据瞬时特征独立决策，容易导致相邻时刻权重出现大幅跳变。引入 \(\boldsymbol{\omega}_{t-1}\) 后，智能体在进行当前决策时能够保留必要的历史上下文，这也为后续增量式动作设计提供了基础。

从方法层看，PA-DQN 的状态空间重构主要体现为两个变化。其一，状态输入由 RL-IFA 中的单一后验误差，扩展为传播、统计、几何和历史决策四类信息的联合描述；其二，状态的作用由“反映当前结果好坏”转为“支撑当前权重应如何调整”。前者决定了智能体能够获得更丰富的环境信息，后者决定了这些信息可以直接服务于后续策略学习。这样处理后，权重分配不再只是对已发生误差的被动修正，而更接近于依据当前环境特征进行前置调节。

需要说明的是，本节给出的状态空间并未改变位置级动态加权融合模型本身，变化仅发生在决策层。RSSI、AoA 和 PDR 三路位置估计仍按既定方式输出初始定位结果，PA-DQN 只是利用式~\eqref{eq:pa_dqn_state} 所构造的状态向量对当前环境进行表征，并在此基础上学习更合适的权重更新策略。后续将在此基础上进一步给出 PA-DQN 的动作空间设计与奖励函数构造。

\begin{figure}[htbp]
    \centering
    % \includegraphics[width=0.82\textwidth]{figures/ch3_state_space_architecture.pdf}
    \vspace{4.5cm}
    \caption{PA-DQN 物理感知状态空间及其与融合决策关系示意图}
    \label{fig:pa_dqn_state_space}
\end{figure}

\subsection{抑制轨迹锯齿的增量式动作空间设计}
\label{subsec:3_4_2}

为了保证输出轨迹的平滑性， PA-DQN 不再采用简单的离散切换逻辑，而是引入增量式动作空间。其基本思想为：当前动作并不直接生成一组全新权重，而是在上一时刻权重基础上施加有限增量更新，即
\begin{equation}
\omega_t = \omega_{t-1} + \Delta \omega_t,
\label{eq:3_14_incremental_weight}
\end{equation}
其中 \(\Delta \omega_t\) 由当前动作决定。更新后再通过归一化约束保证各权重之和为 1。

\subsection{兼顾短时抗扰与长时纠偏的奖励函数设计}
\label{subsec:3_4_3}

本文设计双重奖励函数。设即时误差为 \(\varepsilon_t\)，则基础精度奖励可写为
\begin{equation}
r_t^{(1)}=-\varepsilon_t.
\label{eq:3_15_reward_error}
\end{equation}
为抑制权重突变带来的轨迹锯齿，进一步引入平滑惩罚项：
\begin{equation}
r_t^{(2)}=-\lambda \left\|\omega_t-\omega_{t-1}\right\|_2,
\label{eq:3_16_reward_smooth}
\end{equation}
其中 \(\lambda\) 为平滑惩罚系数。最终奖励函数定义为
\begin{equation}
r_t = r_t^{(1)} + r_t^{(2)}.
\label{eq:3_17_reward_total}
\end{equation}

\subsection{物理感知深度 Q 网络结构设计}
\label{subsec:3_4_4}

在 RL-IFA 基线算法中，状态空间与动作空间均被离散化处理，因而可以采用表格式 Q-Learning 对状态--动作值函数进行迭代更新。然而，随着 PA-DQN 中状态空间由单一误差反馈扩展为包含物理层、统计层、几何层及历史决策信息的高维连续状态向量，传统表格式方法在状态存储规模、泛化能力和训练效率等方面均面临明显限制。一方面，连续状态变量难以通过有限离散网格进行精细刻画，若强行离散化，容易导致状态爆炸和信息损失；另一方面，复杂城市环境下的状态变化具有显著非线性和时变特征，表格式方法难以有效学习不同环境状态与最优权重调整动作之间的复杂映射关系。因此，本文采用深度 Q 网络（Deep Q-Network, DQN）对状态--动作价值函数进行近似表示，在 RL-IFA 序列决策框架基础上构建物理感知深度 Q 网络模型。

设第 \(t\) 个时刻的物理感知状态向量为
\begin{equation}
\mathbf{S}_t=
\left[
\sigma_{\tau,t},
\ \sigma^2_{\mathrm{RSSI},t},
\ D_{\mathrm{RF\mbox{-}PDR},t},
\ \boldsymbol{\omega}_{t-1}
\right]^T,
\label{eq:3_4_4_state_input}
\end{equation}
其中，\(\sigma_{\tau,t}\)、\(\sigma^2_{\mathrm{RSSI},t}\)、\(D_{\mathrm{RF\mbox{-}PDR},t}\) 和 \(\boldsymbol{\omega}_{t-1}\) 的定义已在前文给出。PA-DQN 以该状态向量作为网络输入，通过多层非线性映射输出各候选动作对应的 Q 值估计，即
\begin{equation}
\mathbf{Q}(\mathbf{S}_t;\theta)
=
\left[
Q(\mathbf{S}_t,a_1;\theta),
Q(\mathbf{S}_t,a_2;\theta),
\ldots,
Q(\mathbf{S}_t,a_{|\mathcal{A}|};\theta)
\right]^T,
\label{eq:3_4_4_q_output}
\end{equation}
其中，\(\theta\) 表示深度网络参数，\(|\mathcal{A}|\) 为动作空间规模。网络输出的每一个分量对应在当前状态下执行某一候选动作所能获得的长期累计回报估计，智能体据此选择最优动作完成权重更新。

从结构上看，PA-DQN 采用“状态输入层—隐藏特征提取层—动作价值输出层”的基本架构。输入层负责接收当前时刻的物理感知状态向量；隐藏层通过若干全连接映射和非线性激活，对环境特征、观测波动特征及历史决策信息进行联合表征；输出层则直接给出所有动作的 Q 值估计。与表格式 Q-Learning 相比，这种结构不再显式存储每个离散状态下的动作价值，而是通过参数共享方式学习不同状态之间的潜在关联，从而提升对未见状态和复杂环境条件的泛化能力。

PA-DQN 的核心目标仍然是逼近最优状态--动作值函数 \(Q^*(\mathbf{S}_t,a_t)\)。依据 Bellman 最优性原理，其目标值可写为
\begin{equation}
y_t
=
r_t+\gamma \max_{a'}Q(\mathbf{S}_{t+1},a';\theta^-),
\label{eq:3_4_4_target}
\end{equation}
其中，\(r_t\) 为当前时刻即时奖励，\(\gamma\in(0,1)\) 为折扣因子，\(\theta^-\) 为目标网络参数。于是，当前网络参数 \(\theta\) 的优化目标可以表示为
\begin{equation}
L(\theta)
=
\mathbb{E}_{(\mathbf{S}_t,a_t,r_t,\mathbf{S}_{t+1})\sim \mathcal{D}}
\left[
\left(
y_t-Q(\mathbf{S}_t,a_t;\theta)
\right)^2
\right],
\label{eq:3_4_4_loss}
\end{equation}
其中，\(\mathcal{D}\) 表示经验回放池。通过最小化式~\eqref{eq:3_4_4_loss}，网络不断修正当前状态下各动作的价值估计，使得输出 Q 值逐步逼近最优解。

为了提升训练过程的稳定性，本文在 PA-DQN 中引入经验回放机制和目标网络机制。经验回放的作用在于打破连续样本之间的时序相关性：智能体在与环境交互过程中，将每一步形成的四元组样本 \((\mathbf{S}_t,a_t,r_t,\mathbf{S}_{t+1})\) 存入经验池，并在训练时随机采样小批量数据进行参数更新。这样可避免网络过度受局部轨迹样本影响，从而提高训练稳定性。目标网络则用于提供相对平稳的目标值估计，即在若干训练步内固定 \(\theta^-\) 不变，仅周期性地将当前网络参数复制到目标网络，以减弱目标值随训练同步波动带来的不稳定问题。

在动作选择策略上，PA-DQN 延续强化学习中的探索--利用平衡思想，采用 \(\varepsilon\)-greedy 机制进行动作采样。即在训练阶段，以概率 \(\varepsilon\) 从动作集合 \(\mathcal{A}\) 中随机选择动作，以概率 \(1-\varepsilon\) 选择当前 Q 值最大的动作：
\begin{equation}
a_t=
\begin{cases}
\text{random action}, & \text{with probability } \varepsilon,\\
\arg\max\limits_{a\in\mathcal{A}}Q(\mathbf{S}_t,a;\theta), & \text{with probability } 1-\varepsilon.
\end{cases}
\label{eq:3_4_4_egreedy}
\end{equation}
训练初期较大的 \(\varepsilon\) 有助于扩大对动作空间的探索范围，避免过早陷入局部最优；随着训练进行，逐步减小 \(\varepsilon\) 则有助于强化当前已学得的较优策略。

从方法意义上看，PA-DQN 相较于 RL-IFA 的关键提升，并不只体现在“将表格替换为神经网络”这一形式变化上，而在于它使高维连续状态空间能够被有效用于决策学习。通过深度网络的函数逼近能力，环境复杂度、信号波动、几何一致性以及历史控制信息之间的非线性耦合关系得以在统一模型中表达。这使得智能体能够突破“状态离散化—动作查表”的基线机制，更细致地刻画复杂城市环境中不同观测源可靠性的动态变化规律，并在此基础上形成更具针对性的权重调整策略。

需要指出的是，PA-DQN 的网络结构设计服务于本章的单目标连续定位任务，其重点并不在于追求过深或过宽的模型规模，而在于保证状态表征能力、训练稳定性和部署复杂度之间的平衡。对于低成本无人机平台而言，过于复杂的网络结构会增加训练与推理开销，而过于简单的网络又难以充分表达环境状态与最优动作之间的复杂关系。因此，本文采用适度规模的全连接深度网络作为 Q 函数近似器，并通过经验回放、目标网络和 \(\varepsilon\)-greedy 策略共同构成 PA-DQN 的训练与决策机制。图~\ref{fig:pa_dqn_network_structure} 给出了 PA-DQN 智能体网络结构与训练流程示意图。

\begin{figure}[htbp]
    \centering
    % \includegraphics[width=0.84\textwidth]{figures/ch3_pa_dqn_network_structure.pdf}
    \vspace{4.5cm}
    \caption{PA-DQN 智能体网络结构与训练流程示意图}
    \label{fig:pa_dqn_network_structure}
\end{figure}

\section{仿真预训练策略与网络收敛性分析}
\label{sec:3_5}

本文引入基于射线追踪先验知识的仿真预训练策略，以加快学习过程并提升收敛稳定性。

\section{实验设置与评价指标}
\label{sec:3_6}

本文采用均方根误差（RMSE）和 95\% 分位误差（P95）作为核心评价指标。

\section{实验结果与性能分析}
\label{sec:3_7}

实验结果表明，相较于 EKF 和标准 Q-Learning，PA-DQN 的定位精度分别提升了约 70.7\% 和 63.6\%。

\section{本章小结}
\label{sec:3_8}

本章提出了一种基于物理感知深度强化学习的自适应融合定位方法，有效提升了复杂城市环境中的单目标定位性能。
